\chapter{KẾT LUẬN VÀ GIỚI HẠN ĐÃ BIẾT}

\section{Tổng kết thành quả đạt được}

Đề tài \textbf{``Hệ thống quản lý chiến dịch marketing có tích hợp AI (MarketFlow AI)''} (Mã đề tài: 80300) đã hoàn thành trọn vẹn các mục tiêu nghiên cứu, thiết kế, hiện thực hóa phần mềm và triển khai thực tế đặt ra cho học phần AIA331. Hệ thống hiện đang vận hành ổn định trên môi trường Production VPS tại địa chỉ \url{https://marketing.kienhieu.id.vn} và vượt qua toàn bộ các bài kiểm định tự động.

Các thành quả cốt lõi đạt được bao gồm:
\begin{enumerate}[noitemsep,topsep=2pt]
 \item \textbf{Hiện thực hóa hoàn chỉnh hệ thống quản trị tiếp thị trên lược đồ CSDL 14 bảng}: Quản lý xuyên suốt từ nhân sự, hồ sơ thương hiệu (\texttt{brand\_profile}), từ khóa cấm (\texttt{banned\_keywords}), 7 loại kênh truyền thông, chiến dịch (\texttt{campaigns}) và phân bổ ngân sách (\texttt{campaign\_channels}) đến nội dung, lịch xuất bản, chỉ số hiệu quả (nhập tay \& CSV 2 bước) và thống kê KPI (CTR, CVR, CPC, CPA, CPM, \texttt{used\_pct}).
 \item \textbf{Kiến trúc phân quyền RBAC kết hợp kiểm soát quyền sở hữu và chống tự duyệt}: Bảo vệ chặt chẽ phạm vi dữ liệu theo \texttt{owner\_id} cho vai trò \texttt{MARKETER}, đồng thời áp dụng quy tắc \texttt{CANNOT\_REVIEW\_OWN} ngăn chặn người dùng (kể cả \texttt{MANAGER}) tự phê duyệt bài viết do chính mình soạn thảo.
 \item \textbf{Máy trạng thái duyệt nội dung HITL 5 trạng thái}: Chu trình \texttt{DRAFT} $\rightarrow$ \texttt{PENDING\_REVIEW} $\rightarrow$ \texttt{APPROVED}/\texttt{REJECTED} $\rightarrow$ \texttt{PUBLISHED} tích hợp tự động quét từ khóa cấm khi gửi duyệt, tự động thu hồi phê duyệt khi sửa bài và khóa xóa đối với nội dung đã duyệt (\texttt{CONTENT\_LOCKED}).
 \item \textbf{Kiến trúc AI Tác tử có kiểm soát qua \texttt{ai-runner} và FastMCP (\texttt{/mcp})}: Cô lập tiến trình \texttt{opencode run} dưới tài khoản phi đặc quyền \texttt{mfai}, xác thực qua \texttt{Run-Token} JWT ngắn hạn (10 phút), cung cấp 14 công cụ MCP (9 Read + 5 Propose) và bắt buộc mọi thao tác ghi của AI đi qua thẻ đề xuất \texttt{ai\_actions} (\texttt{PROPOSED}) chờ con người bấm \textit{Xác nhận}.
 \item \textbf{Độ tin cậy đa nhà cung cấp (13 mô hình LLM) và Đánh giá định lượng minh bạch}: Hỗ trợ 13 mô hình trên 5 nền tảng với cơ chế tự động chuyển đổi dự phòng; đánh giá thực nghiệm 4 giai đoạn tối ưu hóa System Prompt đạt tỷ lệ vượt qua 87,5\% (7/8 kịch bản ở bản Production); kiểm thử tự động 151 ca kiểm thử (104 Backend + 18 AI-Runner + 29 Playwright E2E) đạt 100\% Pass.
\end{enumerate}

\section{Đánh giá khách quan các giới hạn kỹ thuật đã biết}

Nhằm duy trì tính trung thực học thuật (\textit{Rule 0}) và định vị đúng ranh giới của giải pháp kỹ thuật hiện tại, nhóm thẳng thắn ghi nhận các điểm giới hạn đã biết (\textit{Known Limitations}) trong Bảng~\ref{tab:known-limitations}.

\begin{table}[!htbp]
\centering
\footnotesize
\renewcommand{\arraystretch}{1.04}
\caption{Đánh giá khách quan các giới hạn kỹ thuật đã biết và phương án nâng cấp}
\label{tab:known-limitations}
\begin{tabularx}{\textwidth}{|L{3.1cm}|Y|Y|}
\hline
\thead{Hạng mục} & \thead{Hiện trạng và giới hạn kỹ thuật} & \thead{Tác động và giải pháp trong tương lai}\\
\hline
Xử lý tuần tự tại \texttt{ai-runner} (\texttt{RUNNER\_}\allowbreak\texttt{CONCURRENCY=1}) & Do mỗi tiến trình \texttt{opencode run} tiêu thụ $\sim$490 MB RSS trên VPS 2,5 GB RAM, dịch vụ \texttt{ai-runner} được giới hạn xử lý tuần tự từng lượt gọi AI qua hàng đợi FIFO. & Khi có nhiều người dùng gửi tin nhắn AI cùng lúc, các yêu cầu sau phải chờ trong hàng đợi. Giải pháp: nâng cấp RAM máy chủ hoặc chuyển sang gọi trực tiếp SDK nhẹ khi mở rộng quy mô lớn.\\
\hline
Độ trễ của các mô hình miễn phí khi gọi chuỗi công cụ MCP & Một lượt hội thoại gọi liên tiếp 2--3 công cụ MCP qua các cổng miễn phí mất khoảng 20--110 giây tùy tải của nhà cung cấp. & Đã giảm thiểu bằng luồng SSE hiển thị tiến trình thời gian thực và hỗ trợ chuyển đổi sang các mô hình tốc độ cao trên Groq LPU.\\
\hline
Tích hợp xuất bản mạng xã hội thủ công ở bước cuối & Hệ thống quản lý trọn vẹn quy trình duyệt nội bộ và lập lịch, nhưng bước xuất bản cuối cùng sử dụng xác nhận liên kết (\texttt{/mark-published} kèm \texttt{published\_url}) thay vì tự động đẩy qua Graph API/TikTok API. & Tích hợp OAuth 2.0 và Webhook chính thức của Meta Marketing API / TikTok Business API khi triển khai cho doanh nghiệp có tài khoản đối tác đã xác minh.\\
\hline
Đa phương thức (Hình ảnh / Video) & Hệ thống hiện tập trung quản lý văn bản tiếp thị (tiêu đề, thân bài, CTA, hashtags) và số liệu định lượng; chưa sinh hoặc chấm điểm trực tiếp tệp ảnh/video quảng cáo. & Mở rộng công cụ MCP kết nối các mô hình thị giác máy tính (Multimodal) để kiểm duyệt chữ trên banner quảng cáo.\\
\hline
\end{tabularx}
\end{table}

\section{Lộ trình phát triển và hoàn thiện tiếp theo}

Từ các giới hạn đã được nhận diện, nhóm xác lập lộ trình kỹ thuật cho các giai đoạn tiếp theo:
\begin{enumerate}[noitemsep,topsep=2pt]
 \item \textbf{Giai đoạn 1 (Tối ưu hóa tốc độ tác tử)}: Kết hợp bộ nhớ đệm ngữ cảnh cho các công cụ đọc cố định (như \texttt{get\_brand\_profile}, \texttt{list\_channels}) nhằm giảm số vòng gọi công cụ MCP trong mỗi lượt hội thoại.
 \item \textbf{Giai đoạn 2 (Tự động hóa xuất bản và đồng bộ số liệu)}: Kết nối trực tiếp API quảng cáo của Meta và Google để tự động đồng bộ chỉ số hiệu quả hàng đêm thay cho thao tác nhập CSV.
 \item \textbf{Giai đoạn 3 (Kiểm duyệt đa phương thức)}: Bổ sung khả năng kiểm tra quy chuẩn thương hiệu trên hình ảnh và video ngắn trước khi phê duyệt.
\end{enumerate}

\section{Lời kết}

Dự án MarketFlow AI là minh chứng thực tế khẳng định trí tuệ nhân tạo tác tử (Agentic AI qua giao thức MCP), khi vận hành trong kiến trúc phân tách an toàn và quy trình \textit{Human-in-the-Loop} (\textit{Proposal-Only}) kiên cố, hoàn toàn có thể trở thành đòn bẩy năng suất thiết thực cho đội ngũ tiếp thị mà vẫn bảo toàn tuyệt đối tính chính xác và trách nhiệm giải trình của dữ liệu doanh nghiệp.
